\documentclass{article}
\usepackage{amsmath}
\usepackage{blindtext}
\usepackage{titlesec}
\usepackage{hyperref}
\begin{document}
\begin{center}
Master Laboratory

Gravitional Constant

Group 29

Performed from March 30th to April 10th 2026

under supervision of

Alexander Doering

Abstract
    
\end{center}
\newpage
\tableofcontents
\newpage
\section{Introduction}

\label{sec: intro}
Gravitation, or gravity, is one of the four fundamental interactions
of nature1, giving rise to attractive forces between objects with mass. Gravitation is the agent
that determines the weight of
objects with mass and causes them to fall when dropped.
Gravitation causes dispersed matter to coalesce, thus accounting
for the existence of planets, stars and galaxies. Gravitation is
responsible for keeping the Earth and the other planets in their
orbits around the Sun; for keeping the Moon
phase locked in its orbit around the Earth; for the formation of tides
in the oceans; and is one
cause of heating of the interiors during the formation of stars and
planets. In a different context,
gravitation is responsible for natural convection in fluids and gases,
causing flow to occur under
the influence of a density gradient.
In 1684, Isaac Newton (4 January 1643- 31 March 1727) formulated the
law of universal gravitation, according to which the mutual attractive
force between two masses, $m_1$ and $m_2$, is proportional to the
product of the masses and inversely proportional to the square of the
distance $r$
between their respective centre of mass (CM).
\begin{equation}
F =\frac{G m_1 m_2}{r^2}
\end{equation}
% The factor of proportionality G, the universal gravitational constant3, is an empirical physical
% constant that appears as in Eq.4.1. It is universal in the sense that it is valid in the whole of
% universe. The gravitational constant is also known as Newton’s constant, and colloquially ’Big G’.
% It must not be confused with ’little g’, g, which, which is not a constant, and has the dimension
% acceleration; it includes other effects, not only the strength of the local gravitational field4. ’Little
% g’ is equivalent to the free-fall acceleration, especially that at Earth’s surface.
% Although Newton’s theory has been superseded, most modern non-relativistic gravitational calcu
% lations are still using some variety of Eq.4.1, because of its simplicity, and sufficient accuracy for
% most Earth bound applications.
% A general treatment of gravitation requires Einstein’s theory of general relativity, which describes
\newpage
\section{Theory}
\label{sec: theory}
In the following sections, some information specific to the optical
\newpage
\section{Setup and Procedure}
In the following sections, some information specific to the optical
\newpage
\section{Data Analysis}
i
\newpage
\section{uncertainty treatment}
i
\newpage
\section{assumption}
i
\newpage
\section{comparison to expectations }
i
\newpage
\section{results}
i
\newpage
\section{attachments}
i
\end{document}